\chapter{Logarithms as Structure-Preserving Transformations}

In the preceding chapters, invariance under change of coordinates was identified
as the defining characteristic of geometric and tensorial objects.
A vector, a bilinear form, a tensor --- each was defined not by its numerical
components, but by the structural relations it preserves under transformation.
The present chapter observes that this same principle operates outside the
linear setting.
The logarithm is an isomorphism between algebraic structures, converting
multiplication into addition while preserving every relation that matters.
The invariant is not a coordinate, a length, or a tensor component ---
it is the algebraic correspondence itself.

% ------------------------------------------------------------------
%\section*{The Central Idea: A Structure-Preserving Map}
% ------------------------------------------------------------------

\medskip

At the most basic level, a logarithm answers the question:
\emph{what exponent produces a given number?}
Yet historically and mathematically,
logarithms owe their power to something deeper:
they convert multiplication into addition --- and they do so exactly.

\medskip

In Chapter~I, an isomorphism was defined as a structure-preserving mapping
between algebraic systems.
The logarithm provides a precise, global realization of that idea.

\medskip

The mapping
\[
\log:\mathbb{R}^+\to\mathbb{R}
\]
is not merely a bijection; it is an isomorphism between two group structures:

\begin{enumerate}
    \item A multiplicative group $(\mathbb{R}^+,\cdot)$.
    \item An additive group $(\mathbb{R},+)$.
    \item An isomorphism
    \[
    \log:(\mathbb{R}^+,\cdot)\longrightarrow(\mathbb{R},+),
    \]
    whose inverse is $\exp:(\mathbb{R},+)\to(\mathbb{R}^+,\cdot)$.
\end{enumerate}

\medskip

The defining property is:
\[
\log(ab)=\log(a)+\log(b),
\qquad
\log(1)=0.
\]

Multiplication corresponds to addition; multiplicative identity corresponds to
additive identity; multiplicative inverses correspond to additive inverses.
No algebraic relation is lost --- only its appearance changes.

\medskip

This is not a numerical trick or a computational shortcut.
It is a structural equivalence: the two groups are, from an algebraic standpoint,
the same object wearing different notation.

% ------------------------------------------------------------------
%\section*{Physical Manifestation: The Slide Rule}
% ------------------------------------------------------------------

\medskip

This isomorphism is made physical in the slide rule.
Distances on a logarithmic scale represent logarithmic values.
Adding distances corresponds to multiplying numbers.

\medskip

The construction of logarithmic scales proceeds by mapping
\[
x \mapsto \log_{10}(x),
\]
placing numbers at distances proportional to their logarithms.
Equal ratios correspond to equal distances.
The structural identity $\log(ab)=\log(a)+\log(b)$ becomes a mechanical operation:
slide one rule against another, and multiplication is performed by addition of lengths.

\medskip

Lengths are not numbers, but they obey the same structural rules.
By exploiting this correspondence, the slide rule becomes an analog computer ---
a physical model of an algebraic isomorphism.

% ------------------------------------------------------------------
%\section*{Logarithms as Area: The Geometric Interpretation}
% ------------------------------------------------------------------

\medskip

Logarithms also admit a purely geometric definition independent of exponentiation.
By construction,
\[
\ln(a)=\int_1^a \frac{1}{x}\,dx,
\]
the signed area under the curve $y=\frac{1}{x}$ from $1$ to $a$.

\medskip

This integral definition makes the isomorphism property transparent.
For $a,b>0$,
\[
\ln(ab)
=\int_1^{ab}\frac{1}{x}\,dx
=\int_1^{a}\frac{1}{x}\,dx+\int_a^{ab}\frac{1}{x}\,dx.
\]
The substitution $x=at$ in the second integral gives
\[
\int_a^{ab}\frac{1}{x}\,dx=\int_1^{b}\frac{1}{t}\,dt=\ln(b),
\]
and therefore $\ln(ab)=\ln(a)+\ln(b)$.
The group isomorphism is a theorem, not an assumption.

\medskip

The integral definition also yields the fundamental derivative formula immediately:
\[
\frac{d}{da}\ln(a)=\frac{1}{a}.
\]

% ------------------------------------------------------------------
%\section*{Rate of Change and Linear Approximation}
% ------------------------------------------------------------------

\medskip

The behavior of $\frac{d}{da}\ln(a) = \frac{1}{a}$ connects the logarithm to
the general framework of derivatives and local linearization.

\medskip

Recall that for a function $y = f(x)$, the derivative at a point $a$ is
\[
f'(a) = \lim_{h\to0}\frac{f(a+h)-f(a)}{h},
\]
the instantaneous rate of change, and geometrically the slope of the tangent line.
Near $a$, a differentiable function behaves approximately like its tangent line:
\[
f(x) \approx f(a) + f'(a)(x-a).
\]

\medskip

This local linearity mirrors the passage from tensors to their components
discussed in Chapter~II: global nonlinear structure becomes tractable when
approximated by its best linear model at each point.
The derivative replaces a nonlinear map by a linear one locally;
the logarithm replaces a multiplicative structure by an additive one globally.
Both are instances of the same organizing principle.

% ------------------------------------------------------------------
%\section*{Numerical Computation: Newton--Raphson Applied to the Logarithm}
% ------------------------------------------------------------------

\medskip

The derivative of the logarithm also enables its efficient numerical computation
via the Newton--Raphson method.
Suppose we wish to find $\ln(x)$, i.e., to solve the equation
\[
e^y = x,
\]
or equivalently $f(y) = e^y - x = 0$.

\medskip

The Newton--Raphson iteration replaces the nonlinear equation by a sequence of
linear ones.
Starting from a guess $y_n$, we approximate $f$ by its tangent line and solve:
\[
0 = f(y_n) + f'(y_n)(y_{n+1}-y_n).
\]
Since $f'(y) = e^y$, this yields
\[
y_{n+1} = y_n - \frac{e^{y_n}-x}{e^{y_n}}
         = y_n + \frac{x - e^{y_n}}{e^{y_n}}.
\]

Starting from $y_0 = 1.5$, successive iterations converge rapidly to
$\ln(5) \approx 1.60944$.

\medskip

The simplicity of this iteration is not accidental.
Working in base $e$ ensures that $\frac{d}{dy}e^y = e^y$,
minimizing algebraic overhead and improving numerical stability.
Logarithms in other bases are recovered by the change-of-base formula:
\[
\log_{b}(x) = \frac{\ln(x)}{\ln(b)}.
\]

% ------------------------------------------------------------------
%\section*{Historical Algorithm: Briggs and Binary Accumulation}
% ------------------------------------------------------------------

\medskip

The same structural idea underlies historical algorithms for constructing
logarithm tables.
Henry Briggs computed base-10 tables by accumulating small multiplicative
corrections; a closely related method suited to binary arithmetic is described
in~\cite{knuth}.

\medskip

The algorithm builds a target number by successive multiplications of the form
\[
\frac{10^k}{10^k-1},
\]
accumulating the corresponding logarithms whenever the partial product does not
overshoot the target.
What appears as computation is, in fact, structure preservation:
each multiplicative step corresponds to an additive increment on the logarithmic
side, faithfully enacting the isomorphism $\log(ab) = \log(a) + \log(b)$.

% ------------------------------------------------------------------
%\section*{The Constant $e$ and Its Representations}
% ------------------------------------------------------------------

\medskip

The natural base $e$ is not an arbitrary choice.
It is the unique base for which $\frac{d}{dx}e^x = e^x$,
and the unique solution to $\int_1^e \frac{1}{x}\,dx = 1$.
Its arithmetic character is revealed by several independent constructions.

\medskip

One classical definition arises from compound interest:
\[
e = \lim_{n\to\infty}\left(1+\frac{1}{n}\right)^n.
\]

Another is provided by its continued fraction expansion:
\[
e = [2;\,1,2,1,1,4,1,1,6,\dots],
\]
written in expanded form as
\[
e
=
2 + \cfrac{1}{1 + \cfrac{1}{2 + \cfrac{1}{1 + \cfrac{1}{1 + \cfrac{1}{4 + \cdots}}}}}.
\]

These representations converge rapidly and are well suited for computation,
but their deeper significance is structural:
$e$ is the natural unit of the isomorphism between $(\mathbb{R}^+,\cdot)$
and $(\mathbb{R},+)$.

% ------------------------------------------------------------------
%\section*{Two Complementary Linearizations}
% ------------------------------------------------------------------

\medskip

The logarithm transforms nonlinear growth into linear motion.
It converts multiplication into addition, products into sums,
and exponential scales into linear ones.

\medskip

Two complementary manifestations of the invariant spirit are therefore visible.
The logarithm is a \emph{global} isomorphism:
it establishes a perfect structural equivalence between two algebraic worlds,
valid everywhere on $\mathbb{R}^+$.
The derivative, by contrast, provides a \emph{local} linearization:
it approximates an arbitrary transformation by its best linear model at each point.

\medskip

In both cases, complexity is rendered intelligible by expressing it through
linear relations.
This is why logarithms appear simultaneously in analysis, geometry, numerical
methods, and physical instruments.
They do not belong to a single domain.
They preserve structure across domains.